\chapter{Egorov, Vitali and Weak-Lp Ideas}
\label{ch:iii08}

Egorov, Vitali, and weak-$L^p$ theory compare convergence modes and exceptional-set control. They clarify how almost-everywhere convergence can become nearly uniform and how tail bounds encode borderline integrability.

\section*{Learning goals}
After this chapter, the reader should be able to:
\begin{itemize}
\item apply Egorov's theorem on finite-measure sets and track the exceptional-set measure explicitly;
\item apply the Vitali convergence principle using uniform integrability and convergence in measure;
\item compute weak-\(L^p\) distribution-function bounds for explicit functions;
\item distinguish strong \(L^p\) membership from weak-\(L^p\) membership by examples;
\item translate decay of level-set measure into weak-norm estimates and back again;
\item construct examples showing why finite measure, covering hypotheses, or strong-integrability assumptions cannot be dropped silently;
\end{itemize}
\section*{Conceptual roadmap}
\[
\boxed{\text{structure}\;\longrightarrow\;\text{estimate}\;\longrightarrow\;\text{limit or representation}\;\longrightarrow\;\text{application}.}
\]

\paragraph{Mathematical setting.}
Functions are measurable and finite a.e.; in Egorov a finite-valued representative of the limit is used. Here $0<p<\infty$ in weak-$L^p$ and Chebyshev formulas. Put $Q_p(f)=\sup_{t>0}t^p\mu(|f|>t)$ and $\|f\|_{p,\infty}=Q_p(f)^{1/p}$. The latter is a quasi-norm: the inclusion $\{|f+g|>t\}\subset\{|f|>t/2\}\cup\{|g|>t/2\}$ yields a quasi-triangle bound. $Q_p$ itself is homogeneous of degree $p$, not one. For Vitali on $\mu(X)<\infty$, uniform integrability means $\sup_f\|f\|_1<\infty$ and $\lim_{\delta\downarrow0}\sup_f\sup_{\mu(E)<\delta}\int_E|f|=0$. Equivalently the tails $\sup_f\int_{|f|>R}|f|$ tend to zero. Small-set control alone misses unbounded constants on an atom.

\section{Convergence in measure}
A sequence converges in measure when $\mu(|f_n-f|>\varepsilon)\to0$ for every $\varepsilon>0$.

\section{Egorov theorem}
On finite-measure spaces, almost-everywhere convergence becomes uniform after removing a set of arbitrarily small measure.


% BEGIN VOL03-EXPANSION III08-example-01
\begin{example}[Egorov in a concrete sequence]\label{ex:iii08-expand-01}
On \([0,1]\), \(x^n\to0\) almost everywhere. Remove
\(E=(1-\eta,1]\), whose measure is \(\eta\). On the complement,
\[
\sup x^n=(1-\eta)^n\to0.
\]
Thus convergence is uniform outside a set of arbitrarily small measure, exactly
as Egorov predicts.
\end{example}
% END VOL03-EXPANSION III08-example-01
\section{Uniform integrability}
A family in $L^1$ is uniformly integrable when its $L^1$ norms are uniformly bounded and its integrals are uniformly absolutely continuous on small measurable sets, as specified above.

\section{Vitali convergence}
Convergence in measure plus uniform integrability upgrades to $L^1$ convergence under the standard finite-measure hypotheses.

\section{Weak-$L^p$}
Weak-$L^p$ controls tails through $\sup_{t>0}t^p\mu(|f|>t)$ rather than directly integrating $|f|^p$.


% BEGIN VOL03-EXPANSION III08-example-02
\begin{example}[A weak-Lp critical singularity]\label{ex:iii08-expand-02}
On \((0,1)\), let \(f(x)=x^{-1/p}\). For \(t\ge1\),
\[
\mu(|f|>t)=t^{-p}.
\]
Hence \(t^p\mu(|f|>t)=1\), so \(f\) is weak-\(L^p\). But
\[
\int_0^1|f|^pdx=\int_0^1x^{-1}dx=\infty,
\]
so it is not strong \(L^p\).
\end{example}
% END VOL03-EXPANSION III08-example-02
\section{Chebyshev inequality}
Strong $L^p$ bounds imply weak-$L^p$ tail estimates.

\section{Subsequence extraction}
Convergence in measure admits almost-everywhere convergent subsequences under standard assumptions.

\section{Borderline behavior}
Power singularities often lie in weak-$L^p$ at the exact exponent where strong $L^p$ fails.


% BEGIN VOL03-EXPANSION III08-example-03
\begin{example}[Borderline power behavior]\label{ex:iii08-expand-03}
For \(x^{-\alpha}\) on \((0,1)\), strong \(L^p\) requires
\(\alpha p<1\). At the critical relation \(\alpha p=1\), the strong integral
diverges logarithmically while the distribution function decays like
\(t^{-p}\). Weak-\(L^p\) therefore records the exact critical endpoint of the
power-law scale.
\end{example}
% END VOL03-EXPANSION III08-example-03
\section{Core structural results}

\begin{theorem}[Egorov theorem]\label{thm:iii08-01}
If $\mu(X)<\infty$ and $f_n\to f$ a.e., then for every $\varepsilon>0$ there is measurable $E$ with $\mu(E)<\varepsilon$ such that convergence is uniform on $X\setminus E$.
\begin{proof}
For each accuracy level, form decreasing tail-exception sets. Their intersection is null; continuity from above on the finite-measure space yields a late stage with small measure, and a diagonal choice handles all accuracy levels.
\end{proof}
\end{theorem}

\begin{theorem}[Chebyshev inequality]\label{thm:iii08-02}
If $f\in L^p$ and $t>0$, then $\mu(|f|>t)\le t^{-p}\|f\|_p^p$.
\begin{proof}
On the superlevel set, $|f|^p\ge t^p$; integrate and rearrange.
\end{proof}
\end{theorem}

\begin{theorem}[Vitali convergence principle]\label{thm:iii08-03}
On a finite-measure space, convergence in measure plus uniform integrability implies $L^1$ convergence.
\begin{proof}
First extract an a.e.-convergent subsequence using summable exceptional-set measures. Fatou and the uniform $L^1$ bound give $f\in L^1$, and Fatou on each measurable set transfers uniform small-set control to $f$. For $A_n=\{|f_n-f|>\eta\}$, the error is at most $\eta\mu(X)+\int_{A_n}(|f_n|+|f|)$. Convergence in measure and uniform absolute continuity control the second term; then let $\eta\downarrow0$.
\end{proof}
\end{theorem}

\section{Worked examples}

\begin{example}[Egorov needs finite measure]\label{ex:iii08-worked-01}
On \(\mathbb R\), let \(f_n=\mathbf1_{[n,n+1]}\). Then \(f_n\to0\)
pointwise. If \(E\) has arbitrarily small finite measure, infinitely many
intervals \([n,n+1]\) retain points outside \(E\), so
\(\sup_{\mathbb R\setminus E}|f_n|=1\) infinitely often. Thus the finite-measure
hypothesis in Egorov cannot simply be dropped.
\end{example}

\begin{example}[\(L^1\)-boundedness is not uniform integrability]\label{ex:iii08-worked-02}
On \((0,1)\), set \(f_n=n\mathbf1_{(0,1/n)}\). Then
\(\|f_n\|_1=1\), but for every \(\delta>0\) one can choose \(n\) with
\(1/n<\delta\) and take \(E=(0,1/n)\):
\[
\int_E|f_n|=1.
\]
Hence the family is not uniformly integrable despite being \(L^1\)-bounded.
\end{example}

\begin{example}[Uniform integrability repairs convergence]\label{ex:iii08-worked-03}
Vitali's principle says that convergence in measure together with uniform
integrability yields \(L^1\) convergence (under the stated finite-measure or
tightness hypotheses). The moving-spike family fails exactly at uniform
integrability, explaining why its pointwise convergence does not become
\(L^1\) convergence.
\end{example}

\begin{example}[A weak-\(L^p\) model singularity]\label{ex:iii08-worked-04}
On \((0,1)\), let \(f(x)=x^{-1/p}\). Then for \(\lambda\ge1\),
\[
m\{|f|>\lambda\}=m(0,\lambda^{-p})=\lambda^{-p}.
\]
Thus \(\sup_{\lambda>0}\lambda\,m\{|f|>\lambda\}^{1/p}<\infty\), so
\(f\in L^{p,\infty}\), while \(|f|^p=x^{-1}\) is not integrable.
\end{example}

\section{Solved dossiers}
Each dossier combines several ideas from the chapter in a longer argument. The aim is to make hypotheses, calculations, and proof strategy explicit.

\begin{problem}[Egorov needs finite measure]\label{prob:iii08-dossier-01}
Give a pointwise-convergent sequence on an infinite-measure space for which the Egorov conclusion fails.
\end{problem}
\begin{solution}
Use \(f_n=\mathbf1_{[n,n+1]}\) on \(\mathbb R\). The supports escape to
infinity, so pointwise convergence holds, but deleting a small finite-measure
set cannot make the convergence uniform.
\end{solution}

\begin{problem}[\(L^1\)-boundedness is not uniform integrability]\label{prob:iii08-dossier-02}
Show that boundedness in \(L^1\) does not imply uniform integrability.
\end{problem}
\begin{solution}
The spike family \(n\mathbf1_{(0,1/n)}\) has unit \(L^1\) norm but
concentrates all mass on sets of arbitrarily small measure.
\end{solution}

\begin{problem}[Uniform integrability repairs convergence]\label{prob:iii08-dossier-03}
Explain which missing property prevents the moving spikes from converging in \(L^1\).
\end{problem}
\begin{solution}
They are not uniformly integrable: unit mass concentrates on sets whose
measure tends to zero.
\end{solution}

\begin{problem}[A weak-\(L^p\) model singularity]\label{prob:iii08-dossier-04}
Give a function in weak \(L^p(0,1)\) that is not in strong \(L^p(0,1)\).
\end{problem}
\begin{solution}
Take \(f(x)=x^{-1/p}\). Its distribution function is of order
\(\lambda^{-p}\), while \(\int_0^1|f|^p=\int_0^1x^{-1}dx=\infty\).
\end{solution}

\begin{problem}[Almost-uniform convergence]\label{prob:iii08-dossier-05}
State Egorov's conclusion on a finite-measure space.
\end{problem}
\begin{solution}
For a.e. pointwise convergence and every \(\varepsilon>0\), there is a
measurable \(E\) with \(\mu(E)<\varepsilon\) such that convergence is uniform
on \(X\setminus E\).
\end{solution}

\begin{problem}[Uniform-integrability tail test]\label{prob:iii08-dossier-06}
Give a common tail criterion for uniform integrability.
\end{problem}
\begin{solution}
A typical criterion is
\(\sup_{f\in\mathcal F}\int_{\{|f|>M\}}|f|\to0\) as \(M\to\infty\), together
with the appropriate finite-measure setting.
\end{solution}

\begin{problem}[Weak-\(L^p\) distribution estimate]\label{prob:iii08-dossier-07}
If \(f\in L^p\), derive a weak-\(L^p\) estimate using Chebyshev.
\end{problem}
\begin{solution}
Chebyshev gives
\(\mu\{|f|>\lambda\}\le\lambda^{-p}\|f\|_p^p\), hence
\(\lambda\mu\{|f|>\lambda\}^{1/p}\le\|f\|_p\).
\end{solution}

\begin{problem}[Convergence in measure from \(L^1\)]\label{prob:iii08-dossier-08}
Show \(L^1\) convergence implies convergence in measure.
\end{problem}
\begin{solution}
Markov's inequality gives
\(\mu(|f_n-f|>\varepsilon)\le\varepsilon^{-1}\|f_n-f\|_1\to0\).
\end{solution}

\begin{problem}[Egorov theorem proof dossier]\label{prob:iii08-dossier-09}
Prove the following structural statement and identify the step where its main hypothesis is used: If $\mu(X)<\infty$ and $f_n\to f$ a.e., then for every $\varepsilon>0$ there is measurable $E$ with $\mu(E)<\varepsilon$ such that convergence is uniform on $X\setminus E$.
\end{problem}
\begin{solution}
For each accuracy level, form decreasing tail-exception sets. Their intersection is null; continuity from above on the finite-measure space yields a late stage with small measure, and a diagonal choice handles all accuracy levels. This proof also explains why the stated hypothesis is structural rather than cosmetic.
\end{solution}

\begin{problem}[Chebyshev inequality proof dossier]\label{prob:iii08-dossier-10}
Prove the following structural statement and identify the step where its main hypothesis is used: If $f\in L^p$ and $t>0$, then $\mu(|f|>t)\le t^{-p}\|f\|_p^p$.
\end{problem}
\begin{solution}
On the superlevel set, $|f|^p\ge t^p$; integrate and rearrange. This proof also explains why the stated hypothesis is structural rather than cosmetic.
\end{solution}

\begin{problem}[Vitali convergence principle proof dossier]\label{prob:iii08-dossier-11}
Prove the following structural statement and identify the step where its main hypothesis is used: On a finite-measure space, convergence in measure plus uniform integrability implies $L^1$ convergence.
\end{problem}
\begin{solution}
First extract an a.e.-convergent subsequence using summable exceptional-set measures. Fatou and the uniform $L^1$ bound give $f\in L^1$, and Fatou on each measurable set transfers uniform small-set control to $f$. For $A_n=\{|f_n-f|>\eta\}$, the error is at most $\eta\mu(X)+\int_{A_n}(|f_n|+|f|)$. Convergence in measure and uniform absolute continuity control the second term; then let $\eta\downarrow0$. This proof also explains why the stated hypothesis is structural rather than cosmetic.
\end{solution}

\begin{problem}[Chapter synthesis]\label{prob:iii08-dossier-12}
Connect the main definitions and structural theorems of this chapter into one proof strategy. Explain what object is controlled, what estimate or closure principle is used, and what conclusion becomes available.
\end{problem}
\begin{solution}
Egorov, Vitali, and weak-$L^p$ theory compare convergence modes and exceptional-set control. They clarify how almost-everywhere convergence can become nearly uniform and how tail bounds encode borderline integrability. The recurring strategy is to encode the object in the chapter's natural measurable, normed, Fourier, or distributional structure, establish the controlling estimate, and then pass to the desired limit or representation.
\end{solution}
\section{Exercises with complete solutions}
The shorter exercise layer remains separate from the solved dossiers.

\begin{exercise}\label{exr:iii08-01}
State and apply the central principle behind Convergence in measure. Give one immediate consequence in the setting of this chapter.
\end{exercise}
\begin{hint}
Egorov needs finite measure; choose the exceptional set first and make uniform convergence hold on its complement.
\end{hint}
\begin{solution}
A sequence converges in measure when $\mu(|f_n-f|>\varepsilon)\to0$ for every $\varepsilon>0$. As an immediate consequence, the same principle can be inserted into later convergence, approximation, transform, or weak-form arguments whenever its hypotheses are verified.
\end{solution}

\begin{exercise}\label{exr:iii08-02}
State and apply the central principle behind Egorov theorem. Give one immediate consequence in the setting of this chapter.
\end{exercise}
\begin{hint}
Quantify the exceptional set by summing the measures of the bad sets in the proof.
\end{hint}
\begin{solution}
On finite-measure spaces, almost-everywhere convergence becomes uniform after removing a set of arbitrarily small measure. As an immediate consequence, the same principle can be inserted into later convergence, approximation, transform, or weak-form arguments whenever its hypotheses are verified.
\end{solution}

\begin{exercise}\label{exr:iii08-03}
State and apply the central principle behind Uniform integrability. Give one immediate consequence in the setting of this chapter.
\end{exercise}
\begin{hint}
For a Vitali covering argument, extract a disjoint subfamily and compare the original set with an expanded union of those balls.
\end{hint}
\begin{solution}
A family in $L^1$ is uniformly integrable when its $L^1$ norms are uniformly bounded and its integrals are uniformly absolutely continuous on small measurable sets, as specified above. As an immediate consequence, the same principle can be inserted into later convergence, approximation, transform, or weak-form arguments whenever its hypotheses are verified.
\end{solution}

\begin{exercise}\label{exr:iii08-04}
State and apply the central principle behind Vitali convergence. Give one immediate consequence in the setting of this chapter.
\end{exercise}
\begin{hint}
Compute the level set \(\{|f|>\lambda\}\) explicitly before evaluating the weak-\(L^p\) quantity.
\end{hint}
\begin{solution}
Convergence in measure plus uniform integrability upgrades to $L^1$ convergence under the standard finite-measure hypotheses. As an immediate consequence, the same principle can be inserted into later convergence, approximation, transform, or weak-form arguments whenever its hypotheses are verified.
\end{solution}

\begin{exercise}\label{exr:iii08-05}
State and apply the central principle behind Weak-$L^p$. Give one immediate consequence in the setting of this chapter.
\end{exercise}
\begin{hint}
Weak-\(L^p\) control is \(\sup_{\lambda>0}\lambda^p\mu(|f|>\lambda)\); optimize in \(\lambda\) only after deriving the level-set measure.
\end{hint}
\begin{solution}
Weak-$L^p$ controls tails through $\sup_{t>0}t^p\mu(|f|>t)$ rather than directly integrating $|f|^p$. As an immediate consequence, the same principle can be inserted into later convergence, approximation, transform, or weak-form arguments whenever its hypotheses are verified.
\end{solution}

\begin{exercise}\label{exr:iii08-06}
State and apply the central principle behind Chebyshev inequality. Give one immediate consequence in the setting of this chapter.
\end{exercise}
\begin{hint}
To separate weak and strong \(L^p\), use the borderline power \(|x|^{-n/p}\) on a suitable domain.
\end{hint}
\begin{solution}
Strong $L^p$ bounds imply weak-$L^p$ tail estimates. As an immediate consequence, the same principle can be inserted into later convergence, approximation, transform, or weak-form arguments whenever its hypotheses are verified.
\end{solution}

\begin{exercise}\label{exr:iii08-07}
State and apply the central principle behind Subsequence extraction. Give one immediate consequence in the setting of this chapter.
\end{exercise}
\begin{hint}
Strong \(L^p\) implies weak \(L^p\) by Chebyshev's inequality; write that one-line estimate first.
\end{hint}
\begin{solution}
Convergence in measure admits almost-everywhere convergent subsequences under standard assumptions. As an immediate consequence, the same principle can be inserted into later convergence, approximation, transform, or weak-form arguments whenever its hypotheses are verified.
\end{solution}

\begin{exercise}\label{exr:iii08-08}
State and apply the central principle behind Borderline behavior. Give one immediate consequence in the setting of this chapter.
\end{exercise}
\begin{hint}
Do not use Egorov on an infinite-measure space without first restricting to a finite-measure subset.
\end{hint}
\begin{solution}
Power singularities often lie in weak-$L^p$ at the exact exponent where strong $L^p$ fails. As an immediate consequence, the same principle can be inserted into later convergence, approximation, transform, or weak-form arguments whenever its hypotheses are verified.
\end{solution}


% BEGIN VOL03-EXPANSION III08-exercises-01
\section{Graded supplementary exercises}
These exercises extend the preserved foundational set and are graded by mathematical role.

\subsection*{Standard computations}

\begin{exercise}[Chebyshev]\label{exr:iii08-expand-01}
If ||f||\_p=3, bound the measure of |f|>t.
\end{exercise}
\begin{hint}
Use Chebyshev.
\end{hint}
\begin{solution}
At most 3\textasciicircum{}p/t\textasciicircum{}p.
\end{solution}

\begin{exercise}[Weak exponent]\label{exr:iii08-expand-02}
If a distribution tail is bounded by 4 t\textasciicircum{}(-2), which weak space is indicated?
\end{exercise}
\begin{hint}
Match the exponent.
\end{hint}
\begin{solution}
Weak L2.
\end{solution}

\begin{exercise}[Critical example]\label{exr:iii08-expand-03}
Is x\textasciicircum{}(-1/2) in L2(0,1)?
\end{exercise}
\begin{hint}
Integrate x\textasciicircum{}(-1).
\end{hint}
\begin{solution}
No, but it is weak L2.
\end{solution}

\begin{exercise}[Egorov hypothesis]\label{exr:iii08-expand-04}
What measure hypothesis does Egorov require?
\end{exercise}
\begin{hint}
Check the theorem.
\end{hint}
\begin{solution}
Finite total measure.
\end{solution}

\begin{exercise}[Measure convergence]\label{exr:iii08-expand-05}
State convergence in measure.
\end{exercise}
\begin{hint}
Use bad sets.
\end{hint}
\begin{solution}
For every epsilon, the measure of |f\_n-f|>epsilon tends to zero.
\end{solution}

\subsection*{Proofs}

\begin{exercise}[Chebyshev proof]\label{exr:iii08-expand-06}
Prove Chebyshev.
\end{exercise}
\begin{hint}
On the superlevel set, |f|\textasciicircum{}p >= t\textasciicircum{}p.
\end{hint}
\begin{solution}
Integrate over that set and rearrange.
\end{solution}

\begin{exercise}[Egorov mechanism]\label{exr:iii08-expand-07}
Explain the exceptional-set construction.
\end{exercise}
\begin{hint}
Track points where the tail is not yet uniformly accurate.
\end{hint}
\begin{solution}
A.e. convergence makes the ultimate bad set null; finite measure plus continuity from above gives small late bad sets.
\end{solution}

\begin{exercise}[Strong to weak]\label{exr:iii08-expand-08}
Prove Lp is contained in weak Lp.
\end{exercise}
\begin{hint}
Use Chebyshev.
\end{hint}
\begin{solution}
The weak tail quantity is bounded by the strong Lp norm to the p.
\end{solution}

\begin{exercise}[Vitali split]\label{exr:iii08-expand-09}
Explain the Vitali error split.
\end{exercise}
\begin{hint}
Separate small pointwise error from a small exceptional set.
\end{hint}
\begin{solution}
Finite measure controls the small-error region; uniform integrability controls the exceptional region.
\end{solution}

\subsection*{Counterexamples and hypothesis testing}

\begin{exercise}[Infinite measure Egorov]\label{exr:iii08-expand-10}
Why can Egorov fail on infinite measure?
\end{exercise}
\begin{hint}
Slow convergence can drift outward.
\end{hint}
\begin{solution}
On $\mathbb R$ with Lebesgue measure, $f_n=1_{[n,\infty)}\to0$ pointwise. Outside any set of finite measure, each tail still contains points, so the supremum of $f_n$ is one for every $n$.
\end{solution}

\begin{exercise}[Weak not strong]\label{exr:iii08-expand-11}
Give a weak-Lp function not in strong Lp.
\end{exercise}
\begin{hint}
Use x\textasciicircum{}(-1/p) on (0,1).
\end{hint}
\begin{solution}
Its tail has critical t\textasciicircum{}(-p) decay while its p-th power integral diverges.
\end{solution}

\begin{exercise}[Measure not L1]\label{exr:iii08-expand-12}
Give convergence in measure without L1 convergence.
\end{exercise}
\begin{hint}
Use moving spikes.
\end{hint}
\begin{solution}
n on (0,1/n) converges in measure to zero but has L1 norm 1.
\end{solution}

\subsection*{Applications and investigations}

\begin{exercise}[Critical PDE tails]\label{exr:iii08-expand-13}
Why do weak spaces appear at critical PDE exponents?
\end{exercise}
\begin{hint}
They preserve scale-invariant tail control when strong integrability fails.
\end{hint}
\begin{solution}
They capture borderline kernels and singular solutions.
\end{solution}

\begin{exercise}[Uniform integrability]\label{exr:iii08-expand-14}
Why is uniform integrability a concentration condition?
\end{exercise}
\begin{hint}
Small sets cannot hold large integral mass uniformly.
\end{hint}
\begin{solution}
It rules out moving spike behavior.
\end{solution}

\subsection*{Challenge problems}

\begin{exercise}[Weak quasi-norm]\label{exr:iii08-expand-15}
Compute the weak-Lp tail quantity for x\textasciicircum{}(-1/p) on (0,1).
\end{exercise}
\begin{hint}
Handle t<1 and t>=1.
\end{hint}
\begin{solution}
The supremum of t\textasciicircum{}p times the superlevel measure is 1.
\end{solution}

\begin{exercise}[A.e. subsequence]\label{exr:iii08-expand-16}
Why can convergence in measure yield an a.e.-convergent subsequence?
\end{exercise}
\begin{hint}
Choose a subsequence with summable bad-set measures.
\end{hint}
\begin{solution}
A Borel-Cantelli/diagonal argument gives eventual accuracy almost everywhere.
\end{solution}
% END VOL03-EXPANSION III08-exercises-01
\section*{Chapter summary}
The chapter now supplies a canonical theorem-and-dossier layer for subsequent measure, Fourier, distributional, Sobolev, and PDE arguments.
